\scanpage{14}

sont vérifiées. (Lorsque $X$ est également lisse, \unsure{\deleted{a}c|\deleted{c}a}'est contenu dans XVI 3.7 ; le cas général - sans même supposer $Y$ plat sur $S$, d'ailleurs - s'en déduit par \deleted{d\unsure{u|e}\unsure{n|s}} dévissage facile, et il est traité d'ailleurs dans SGA 2 XIV 1, conjonction de 1.10 et de 1.18.) La class\unsure{\deleted{\textsuperscript{u}}e|\deleted{e}\textsuperscript{u}|ë} $\handwritten{\phi}_f$ de (4.5) peut alors s'interpréter aussi comme un élément \deleted{\unsure{de|du}}

\[
\deleted{\handwritten{\phi}_f}\handwritten{\gamma_{X/Y}\in}\mathrm{H}^{\handwritten{2}(d-d')}_{\unsure{\inserted{\handwritten{X}}|\deleted{\unsure{X|Y}}\handwritten{X}}}(\unsure{\handwritten{Y}|\handwritten{X}|\deleted{\unsure{X|Y}}\handwritten{Y}},A_{\unsure{\handwritten{Y}|\deleted{\unsure{X|Y}}\handwritten{Y}}}(d-d')).
\tag*{(4.8)}
\]

Cette classe s'appelle la \uline{classe de cohomologie à supports dans} $X$ \uline{de} $X$ \uline{dans} $Y$, et son image

\[
\handwritten{\gamma}_{X/Y}\handwritten{\in}\mathrm{H}^{\handwritten{2}(d-d')}(Y,A_Y(d-d'))
\tag*{(4.9)}
\]

\uline{la classe de cohomologie} (sous-entendue: à supports quelconques) \uline{de} $X$ \uline{dans} $Y$. Le formalisme des classes de cohomologie\inserted{\handwritten{(4.8) (4.9)}} ainsi \deleted{associées}\inserted{\handwritten{obtenues}} \deleted{aux \unsure{e}} (une fois la définition étendue par linéarité à des cycles algébriques relatifs sur $Y$) sera étudié systématiquement dans SGA 5 IV. Notre propos ici est d\unsure{è|é|e|\deleted{'}e|\deleted{t}e|\deleted{e}'|\deleted{e}t}\deleted{en} donner \deleted{certaines pr} les propri\unsure{\deleted{t}é|\deleted{é}t}tés les plus élémentaires de la classe locale (4.6), en attachant notre attention surtout au cas où $X$ est également lisse sur $S$.

Voici une liste de propriétés à expliciter qui me semble raisonnable\deleted{\unsure{;|:}} (où on pose $\deleted{p}\handwritten{\delta}=d-d'$\deleted{\unsure{;|:}}\handwritten{)}:

a) Si $X$ est lisse sur $S$, et si $G$ est localement\inserted{(ét)} sur $Y$ isomorphe à l'image inverse d'un complexe sur $S$,\inserted{\handwritten{\unsure{p.ex.|p.cx.} $\unsure{G|C}=(A)_Y$,}} alors (4.4) est un isomorphisme. (C'est essentiellement contenu dans XVI 3.7).

b) Si $\deleted{p}\handwritten{\delta}=0$ i.e. $f$ un isomorphisme, (4.4) est l'identité i.e. $\handwritten{\phi}_f=1$ dans (4.5) ou (4.6).

\handwritten{\tikz[baseline=0pt,overlay,x=1cm,y=1cm]{\draw[line width=.45pt] (.07,.69) .. controls (.10,.83) and (-.25,.90) .. (-.43,.60) .. controls (-.65,.32) and (-.66,-.05) .. (-.48,-.31) .. controls (-.30,-.55) and (.09,-.54) .. (.42,-.32) .. controls (.70,-.11) and (.84,.22) .. (.66,.53) .. controls (.57,.77) and (.39,.85) .. (.16,.81);}}c) Si $\deleted{p}\handwritten{\delta}=1$, de sorte que $X$ est un diviseur relatif sur $Y/S$, alors $\handwritten{\phi}_f$ de (4.5) i.e. (4.8), et a fortiori de (4.6) ou (4.9), est donné par la théorie de Kummer comme d'habitude.

d) Compatibi\unsure{\deleted{m}l|\deleted{l}m|\deleted{n}l|\deleted{l}n}ité de $\unsure{\underline{\handwritten{\phi}}|\handwritten{\phi}}_f$ avec \uline{transitivité} et \uline{intersections} (\unsure{\deleted{v}c|\deleted{c}v}es deux formules sont essentiellement identiques), dans le cas \uline{lisse}\unsure{\handwritten{,}|\handwritten{;}|\deleted{;}|\deleted{,}|\deleted{.}}

% archive sheet number: 13
